\chapter{Who Owns the Means of Computation}
\label{sec:26}
Part III asked who, other than the operator, holds what binds the operator. This Part asks who owns the room, and who teaches the people and the machines that stand in it. Arendt would have said the first of those doesn't belong in politics. The French Revolution, on her account, went wrong when the misery of the poor entered the public realm: necessity crowded out freedom, and questions of livelihood belonged to administration \autocite{arendt1963revolution}. The record of who refused points the other way. Whether saying no is survivable turns, for most people, on whether they can afford to lose the job, and Dewey's public forms around consequences of exactly that kind, once the people they fall on find out they are theirs \autocite{dewey1927public}. Ownership is where the price of refusing gets set, so it is a political question, and this Part treats it as one.

The floor's first requirement cuts across ownership: whoever operates a deployment can't hold what binds it, whether the operator is a firm or a state, and Maven and Gaza both had a state for their operator. That calls for separating functions. It doesn't by itself say who should own anything. Ownership enters at three points. Custody has to be paid for by someone other than the operator or the people it protects. Formation, where a system's commitments are set, is under commercial selection when a market does it. And the offer that makes refusing survivable is priced by whoever sells it. The state is the hardest operator. So each of the three is asked two questions: what fails if a market provides it, and what fails if a government holds it. The answers divide the institutions I propose into two halves.
